\documentclass[10pt,a4paper]{article}
\usepackage[margin=0.68in]{geometry}
\usepackage{amsmath}
\usepackage{booktabs}
\usepackage{graphicx}
\usepackage{microtype}
\usepackage{xcolor}
\usepackage{hyperref}
\usepackage{enumitem}
\setlist{nosep,leftmargin=*}
\setlength{\parindent}{0pt}
\setlength{\parskip}{3pt}
\setlength{\textfloatsep}{7pt}
\setlength{\floatsep}{6pt}
\graphicspath{{../outputs/part_a/figures/}}
\hypersetup{colorlinks=true,urlcolor=blue!60!black,linkcolor=black}
\newcommand{\StudentName}{Student}
\InputIfFileExists{student-info.tex}{}{}

\title{SC4001 Programming Assignment 1\\Part A: HDB Resale Price Prediction}
\author{\StudentName}
\date{}

\begin{document}
\maketitle
\vspace{-1.4em}

\section*{Generative-AI declaration}
OpenAI Codex was used for debugging, figure generation and report drafting.
Its output was verified using unit tests, deterministic
reruns, independent aggregation checks, and visual inspection of all figures.

\section{Experimental protocol}
The data were split chronologically: 2017--2020 for training (87,370 rows),
2021 for validation (29,057), and 2022 for testing (26,702). For model
selection, categorical mappings and continuous-feature statistics were fitted
on 2017--2020 only. After selecting hyperparameters and an epoch budget, the
same preprocessing procedure was refitted on the combined 2017--2021 data and
the model was retrained from scratch. The 2022 test curve was logged only
because the task requested it; it was never used for model or epoch selection.

All models used Adam (learning rate $10^{-2}$), batch size 512, at most 20
epochs, and early stopping with patience 4 and minimum RMSE improvement
$10^{-4}$. CPU training used one thread to make Ray trials and local retraining
numerically reproducible. A2 and A3 used seeds
$\{42,123,456,789,2026\}$ and report mean $\pm$ sample standard deviation.

\section{A1: Baseline and hyperparameter search}
The supplied baseline embeds four categorical features, concatenates the
embeddings with six standardised continuous features, and passes the result
through Linear--LayerNorm--ReLU--Linear layers. Its output is converted from a
standardised target back to Singapore dollars (SGD).

\begin{center}
\includegraphics[width=\textwidth]{a1_baseline_block_diagram.png}
\end{center}

With category cardinalities $(13,27,44,18)$ (including an unknown-category
index) and embedding width 8, the four embedding tables contain 816
parameters. The input width is $4\times8+6=38$, so the complete count is
\[
816+(38\times64+64)+(64+64)+(64+1)=\boxed{3505}.
\]

Ray Tune evaluated all $5\times5=25$ combinations of hidden width
$\{16,32,64,128,256\}$ and common embedding dimension
$\{2,4,8,12,16\}$. The lowest best validation RMSE was 59,666.52 SGD at
hidden width 256 and embedding dimension 4; its best epoch was 6. The selected
network has 7,065 trainable parameters.

\begin{center}
\includegraphics[width=0.72\textwidth]{a1_grid_search_heatmap.png}
\end{center}

\subsection*{Learning behaviour and final test result}
\begin{center}
\begin{minipage}[t]{0.49\textwidth}
\centering
\includegraphics[width=\linewidth]{a1_selection_learning_curve.png}
\end{minipage}\hfill
\begin{minipage}[t]{0.49\textwidth}
\centering
\includegraphics[width=\linewidth]{a1_final_learning_curve.png}
\end{minipage}
\end{center}

The training RMSE decreased almost monotonically, whereas validation RMSE was
noisy and reached its minimum at epoch 6 before worsening. The large
train--validation gap indicates overfitting, while the much higher 2022 error
also suggests temporal distribution shift. Thus the model is not globally
underfitting: it has enough capacity to fit the training data, but generalises
less well to later years. Increasing hidden width generally helped, but the
non-monotonic embedding-dimension results show that simply enlarging every
component was not consistently beneficial.

After refitting preprocessing on 2017--2021 and training for the fixed six
epochs, the seed-42 model achieved
\[
\boxed{\mathrm{RMSE}=97{,}732.09\ \mathrm{SGD}},\qquad
\boxed{R^2=0.6705}
\]
on 2022. The test curve fluctuates, but none of its intermediate values was
used to choose the final epoch.

\section{A2: Ablation studies}
Each variant used the A1 hyperparameters, preprocessing, and six-epoch budget.
``Continuous only'' removed all four categorical embeddings; ``Sigmoid'' kept
the original inputs but replaced the hidden ReLU with a sigmoid.

\begin{center}
\includegraphics[width=0.83\textwidth]{a2_ablation_comparison.png}
\end{center}

\begin{center}
\begin{tabular}{lrrrr}
\toprule
Model & RMSE mean & RMSE SD & $R^2$ mean & $R^2$ SD \\
\midrule
Baseline ReLU & 98,059.90 & 3,498.80 & 0.6680 & 0.0240 \\
Continuous only & 111,213.11 & 7,915.25 & 0.5716 & 0.0604 \\
Baseline Sigmoid & 100,416.22 & 4,504.27 & 0.6516 & 0.0311 \\
\bottomrule
\end{tabular}
\end{center}

Removing categorical features increased mean RMSE by 13,153.21 SGD (13.4\%)
and lowered mean $R^2$ by 0.0963, with greater seed sensitivity. Town, flat
model/type, storey range, and month therefore carry useful information not
fully recoverable from the continuous inputs. Replacing ReLU with Sigmoid
caused a smaller degradation of 2,356.33 SGD (2.4\%) and $0.0164$ in $R^2$.
This is consistent with Sigmoid saturation and weaker gradient propagation,
although the experiment alone cannot isolate optimisation from representational
effects.

\section{A3: Wide-and-deep model}
The same concatenated embedding-plus-continuous vector was sent to two branches.
The deep branch retained the A1 MLP, while the wide branch was a single linear
layer. In standardised target space,
\[
\hat y = f_{\mathrm{deep}}(x)+\lambda f_{\mathrm{wide}}(x),
\quad \lambda\in\{0,0.25,0.5,1,2\}.
\]
All A1 hyperparameters were fixed. Seed 42 and the 2017--2020/2021 split were
used only to select $\lambda$; the lowest validation RMSE, 58,076.59 SGD,
occurred at $\lambda=2$.

\begin{center}
\begin{minipage}[t]{0.48\textwidth}
\centering
\includegraphics[width=\linewidth]{a3_lambda_search.png}
\end{minipage}\hfill
\begin{minipage}[t]{0.50\textwidth}
\centering
\includegraphics[width=\linewidth]{a3_model_comparison.png}
\end{minipage}
\end{center}

\begin{center}
\begin{tabular}{lrrrr}
\toprule
Model & RMSE mean & RMSE SD & $R^2$ mean & $R^2$ SD \\
\midrule
Baseline ReLU & 98,059.90 & 3,498.80 & 0.6680 & 0.0240 \\
Wide-and-deep ($\lambda=2$) & 98,714.00 & 5,698.09 & 0.6630 & 0.0398 \\
\bottomrule
\end{tabular}
\end{center}

The wide path can learn direct additive effects without forcing them through
the nonlinear hidden layer, potentially making low-order trends easier to
optimise. Here it improved the single selection run but did not provide a
robust test benefit: across five seeds its mean RMSE was 654.11 SGD (0.67\%)
worse and its variability was higher. A likely limitation is redundancy: the
deep branch can already represent linear effects, while the extra path changes
optimisation and introduces additional parameters. Selecting $\lambda$ with
one seed also confounds architecture quality with seed-specific noise; nested
multi-seed selection would be stronger but was outside the prescribed setup.

\section{A4: Captum FeatureAblation}
Captum FeatureAblation was applied to every 2022 test sample for the seed-42 A1
baseline and A3 wide-and-deep model. Continuous features were replaced by their
2017--2021 training means (zero after standardisation), and categorical
features by their training modes. Importance is the test-set mean absolute
change in predicted price, in SGD.

\begin{center}
\includegraphics[width=0.92\textwidth]{a4_feature_importance.png}
\end{center}

Both models had the same top three features: \texttt{floor\_area\_sqm},
\texttt{flat\_model\_type}, and \texttt{remaining\_lease\_years}. Their mean
absolute attributions were respectively 67,840, 45,104, and 42,465 SGD for the
baseline, and 63,801, 46,078, and 41,891 SGD for wide-and-deep. This agreement
suggests that the added linear path did not fundamentally change the dominant
predictive signals, although it redistributed some attribution among location
features.

These values are not causal effects. Feature ablation can create unrealistic
feature combinations, correlated predictors can share or transfer importance,
and the ranking depends on the chosen mean/mode baselines. Consequently, the
plot supports a comparison of model reliance under this perturbation scheme,
not claims about the causal determinants of resale price.

\end{document}
